\chapter{T-03: McNemar pareado mBERT vs XLM-R}

\section{Hallazgo}

El informe comparativo señaló que comparar dos F1-macro sin prueba pareada es insuficiente. En la fecha del informe faltaba el vector de predicciones mBERT alineado al test.

\section{Implementación}

\begin{enumerate}[noitemsep]
  \item Exportación de predicciones mBERT: \texttt{reports/metrics/mbert\_test\_predictions.csv} ($n=460$).
  \item Vector congelado XLM-R: \texttt{reports/metrics/xlmr\_test\_predictions.csv}.
  \item Prueba McNemar: \texttt{scripts/run\_mcnemar.py} $\rightarrow$ \texttt{mcnemar\_tests.csv}.
\end{enumerate}

\section{Resultado}

\begin{table}[H]
  \centering
  \caption{Pruebas de McNemar (extracto)}
  \label{tab:mcnemar}
  \begin{tabular}{llcccccc}
\toprule
Modelo A & Modelo B & Estadístico & p-valor & b & c & n & Sig. \\
\midrule
mBERT & LogisticRegression & 0.5647 & 0.277999 & 48 & 37 & 460 & No \\
mBERT & LinearSVC & 0.5402 & 0.520292 & 47 & 40 & 460 & No \\
mBERT & MultinomialNB & 0.7063 & 0.000004 & 89 & 37 & 460 & Sí \\
mBERT & XLM-RoBERTa & 0.4098 & 0.200031 & 25 & 36 & 460 & No \\
\bottomrule
\end{tabular}

\end{table}

La comparación mBERT vs XLM-RoBERTa arroja $p=0.200031$ ($b=25$, $c=36$, $n=460$, $\alpha=0.05$): \textbf{no se rechaza la hipótesis nula}. La diferencia puntual de F1-macro no es estadísticamente concluyente con este tamaño de muestra.

Frente a MultinomialNB la superioridad de mBERT sí alcanza significancia ($p < 0.001$).

\section{Interpretación}

Este resultado cierra el hallazgo crítico \#4 de la auditoría. Debe presentarse como \emph{equivalencia estadística}, no como superioridad de mBERT sobre XLM-R.

\textbf{Criterio de cierre:} fila \texttt{xlm-roberta} en \texttt{mcnemar\_tests.csv} y prosa generada en el informe principal.
